% !TEX root = ../Projektdokumentation.tex
\subsection{Projektorganisation und Zeitplanung}
\label{sec:Projektorganisation}
Die Projektdurchführung erfolgte im für die FIAE-Projektarbeit vorgesehenen
Zeitbudget von 80 Stunden. Zur Steuerung des Vorhabens wurde die Arbeit in
Analyse, Soll-Konzept, Entwurf, Implementierung, Test, Inbetriebnahmevorbereitung
und Dokumentation gegliedert. So war von Anfang an klar, wie viel Zeit für
Analyse, Entwicklung, Tests und Dokumentation zur Verfügung steht.

\begin{table}[htbp]
\centering
\begin{tabular}{|p{9cm}|r|}
\hline
\textbf{Projektphase} & \textbf{Stunden} \\
\hline
Analyse der Ausgangssituation und Anforderungsklärung & 8 \\
\hline
Soll-Konzept und Variantenvergleich & 10 \\
\hline
Entwurf von Datenmodell, Rollenmodell und Oberfläche & 12 \\
\hline
Implementierung der Anwendung & 28 \\
\hline
Tests, Fehlerkorrekturen und Abnahmevorbereitung & 10 \\
\hline
Kundendokumentation und Projektdokumentation & 12 \\
\hline
\textbf{Summe} & \textbf{80} \\
\hline
\end{tabular}
\caption{Geplante Zeitaufteilung für das Projekt}
\label{tab:Zeitplanung}
\end{table}

Die Umsetzung erfolgte schrittweise. Nach der Erhebung der Anforderungen
wurden die Teilbereiche Firmen, Benutzer, Mail-Adressen, IP-Adressen,
Authentifizierung und Auditierung getrennt betrachtet und danach umgesetzt.
Rückmeldungen aus dem internen IT-Umfeld konnten dadurch direkt in den Entwurf
und die Implementierung einfließen.

Materiell wurde ein vorhandener Entwicklungsarbeitsplatz mit Python~3.12,
FastAPI, SQLModel, SQLite und einer lokalen Testumgebung genutzt. Externe
Lizenzkosten fielen für die im Projekt eingesetzten Kernkomponenten nicht an.
